% sec:spec-chebyshev — Chebyshev Polynomials and Gauss-Lobatto Nodes
\section{Chebyshev Polynomials and Gauss-Lobatto Nodes}
\label{sec:spec-chebyshev}

The previous section showed that clustered nodes cure the Runge
phenomenon, and that the Chebyshev zeros minimise the node polynomial
over all monic polynomials of a given degree.  But where do these
special nodes come from?  The answer lies in a single change of
variable: the map $x = \cos\theta$ transforms polynomial
approximation on $[-1, 1]$ into trigonometric approximation on
$[0, \pi]$, and the rich theory of Fourier series carries over
wholesale.  The Chebyshev polynomials that emerge from this
correspondence carry algebraic, orthogonality, and minimax properties
that explain why the Chebyshev--Gauss--Lobatto (CGL) nodes of
\cref{eq:cgl-nodes} --- the extrema of these polynomials --- are the
natural choice for spectral collocation.
\coderef{spectral}{rs}

\paragraph{The cosine definition.}

The $n$-th Chebyshev polynomial of the first kind is defined by the
trigonometric identity
%
\begin{equation}
  T_n(x) = \cos(n\,\arccos x) \,,
  \qquad x \in [-1, 1] \,.
  \label{eq:chebyshev-def}
\end{equation}
%
Setting $\theta = \arccos x$ reveals the essential idea: $T_n$ is
simply $\cos(n\theta)$ written in the variable
$x = \cos\theta$.  The first few are $T_0(x) = 1$,
$T_1(x) = x$, and $T_2(x) = 2x^2 - 1$ (the double-angle
formula).  Despite the transcendental appearance of the definition,
$T_n$ is a polynomial of degree exactly $n$ with integer
coefficients and leading coefficient $2^{n-1}$ for $n \geq 1$.
\histnote{Chebyshev introduced these polynomials in the 1850s in the
context of mechanism design --- specifically, approximating a
rectilinear linkage --- long before their role in numerical analysis
was recognised.}

\paragraph{Three-term recurrence.}

The trigonometric identity
$\cos((n+1)\theta) = 2\cos\theta\,\cos(n\theta) - \cos((n-1)\theta)$
translates directly into the polynomial recurrence
%
\begin{equation}
  T_{n+1}(x) = 2x\,T_n(x) - T_{n-1}(x) \,,
  \label{eq:chebyshev-recurrence}
\end{equation}
%
with initial conditions $T_0(x) = 1$, $T_1(x) = x$.  This
recurrence generates every Chebyshev polynomial from the two
simplest, and is the standard method for evaluating $T_n(x)$ at a
given point --- each step costs one multiply and one subtract, with no
trigonometric function calls.  The recurrence also reveals that the
leading coefficient of $T_n$ doubles at each step, confirming the
factor $2^{n-1}$ for $n \geq 1$.
\implnote{The codebase's \texttt{ChebyshevBasis} in
\codemodule{spectral} uses the cosine formula
(\cref{eq:chebyshev-def}) to place the CGL nodes directly, rather
than evaluating the recurrence.  The recurrence appears in the
Legendre polynomial evaluation (\cref{sec:spec-legendre}), where the
analogous three-term relation is the primary computational tool.}

\paragraph{Zeros and extrema.}

The zeros and extrema of $T_n$ are direct consequences of the
cosine definition.  Setting $T_n(x) = \cos(n\theta) = 0$ gives
$n\theta = (2k+1)\pi/2$, so the $n$ zeros are
%
\begin{equation}
  x_k = \cos\!\left(\frac{2k + 1}{2n}\,\pi\right) ,
  \qquad k = 0, 1, \ldots, n-1 \,.
  \label{eq:chebyshev-zeros-Tn}
\end{equation}
%
These are exactly the Chebyshev zeros of $T_n$ --- the same nodes
that appeared in \cref{eq:chebyshev-zeros} of the previous section
with $n = N + 1$.  The extrema of $T_n$ occur where
$\cos(n\theta) = \pm 1$, i.e.\ $n\theta = j\pi$, giving
the $n + 1$ points
%
\begin{equation}
  x_j = \cos\!\left(\frac{j}{n}\,\pi\right) ,
  \qquad j = 0, 1, \ldots, n \,,
  \label{eq:chebyshev-extrema}
\end{equation}
%
which are the CGL nodes of \cref{eq:cgl-nodes} with $n = N$.
At these points, $T_n(x_j) = \cos(j\pi) = (-1)^j$: the polynomial
alternates between $+1$ and $-1$ at its extrema, touching both
bounds at every one of $n + 1$ equally spaced points in
$\theta$-space.

\paragraph{The minimax property.}

The alternation $T_n(x_j) = (-1)^j$ at the extrema is the key to
the minimax optimality claimed in the previous section.  Among all monic
polynomials of degree $n$ on $[-1, 1]$, the scaled Chebyshev
polynomial $\tilde{T}_n(x) = T_n(x) / 2^{n-1}$ has the smallest
possible supremum norm:
%
\begin{equation}
  \max_{x \in [-1,1]} \abs{\tilde{T}_n(x)}
    = \frac{1}{2^{n-1}} \,.
  \label{eq:chebyshev-minimax}
\end{equation}
%
Any other monic polynomial of degree $n$ must have a larger
supremum norm.  The proof is a short contradiction argument using
the equioscillation theorem: suppose a monic polynomial $q_n$
satisfies $\norm{q_n}_\infty < 2^{-(n-1)}$.  Then
$\tilde{T}_n - q_n$ is a polynomial of degree at most $n - 1$
(the leading terms cancel), but
$\tilde{T}_n(x_j) = (-1)^j / 2^{n-1}$ has magnitude exceeding
$\norm{q_n}_\infty$ at each of the $n + 1$ extremal points.
Since $\abs{\tilde{T}_n(x_j)} > \norm{q_n}_\infty \geq \abs{q_n(x_j)}$,
the sign of $\tilde{T}_n(x_j) - q_n(x_j)$ equals the sign of
$\tilde{T}_n(x_j) = (-1)^j$.  The difference $\tilde{T}_n - q_n$
therefore changes sign between each consecutive pair of extremal
points, giving at least $n$ zeros --- impossible for a polynomial of
degree $n - 1$.

The minimax property explains why the Chebyshev zeros produce the
smallest node polynomial: the monic polynomial
$\omega_{N+1}(x) = \prod_{j=0}^{N}(x - x_j)$
of \cref{eq:node-polynomial} is minimised in supremum norm precisely
when it equals $\tilde{T}_{N+1}$, i.e.\ when the nodes are the
zeros of $T_{N+1}$.
\physnote{The equioscillation property is also the reason Chebyshev
nodes suppress the Runge phenomenon: the node polynomial oscillates
\emph{uniformly} across $[-1,1]$, so no region of the interval
suffers disproportionately large interpolation error.  Equidistant
nodes lack this balance --- the node polynomial is orders of magnitude
larger near the endpoints.}

\paragraph{Orthogonality.}

The Chebyshev polynomials are orthogonal with respect to the
weighted inner product
%
\begin{equation}
  \int_{-1}^{1} T_m(x)\,T_n(x)\,
    \frac{\d x}{\sqrt{1 - x^2}}
  = \begin{cases}
      0 & m \neq n \\
      \pi & m = n = 0 \\
      \pi/2 & m = n \geq 1
    \end{cases}
  \label{eq:chebyshev-ortho}
\end{equation}
%
This follows directly from the substitution $x = \cos\theta$: the
weight function $(1 - x^2)^{-1/2}$ cancels the Jacobian of the
change of variable, and the integral reduces to
$\int_0^\pi \cos(m\theta)\,\cos(n\theta)\,\d\theta$, which is a
standard Fourier orthogonality integral.  The weight
$(1 - x^2)^{-1/2}$ diverges at $x = \pm 1$, reflecting the
clustering of the $\theta$-grid near the endpoints of $[-1,1]$
--- the same clustering that suppresses the Runge phenomenon.
\warnnote{The Chebyshev weight $(1-x^2)^{-1/2}$ differs from the
Legendre weight, which is simply $1$.  The two families are
orthogonal with respect to different inner products, and the choice
of weight function determines which nodes (CGL vs LGL) are natural
for a given application.  \Cref{sec:spec-legendre} develops the
Legendre case.}

\paragraph{Chebyshev expansion.}

Because the Chebyshev polynomials are orthogonal
(\cref{eq:chebyshev-ortho}), any function
$f \in L^2_w[-1,1]$ (square-integrable with respect to
the Chebyshev weight) in a Chebyshev series, converging in the
weighted $L^2$ norm:
%
\begin{equation}
  f(x) = \sum_{n=0}^{\infty} a_n\,T_n(x) \,,
  \qquad
  a_n = \frac{2}{\pi\,c_n}
    \int_{-1}^{1} f(x)\,T_n(x)\,
    \frac{\d x}{\sqrt{1 - x^2}} \,,
  \label{eq:chebyshev-expansion}
\end{equation}
%
where $c_0 = 2$ and $c_n = 1$ for $n \geq 1$.  The coefficients
$a_n$ are Fourier cosine coefficients in disguise: under
$x = \cos\theta$, the expansion becomes
$f(\cos\theta) = \sum a_n \cos(n\theta)$, and the $a_n$ are
computed by the standard cosine transform.  For smooth functions,
the coefficients $a_n$ decay faster than any power of $1/n$ ---
this is the spectral convergence that gives spectral methods their
name, developed quantitatively in \cref{sec:spec-convergence}.

\paragraph{Discrete orthogonality and aliasing.}

Continuous Chebyshev expansions are elegant but impractical: the
integral in \cref{eq:chebyshev-expansion} is rarely available in
closed form.  The CGL nodes provide an exact quadrature rule for
polynomials: for any polynomials $T_m$ and $T_n$ with
$m + n \leq 2N - 1$, the Gauss--Lobatto quadrature on the $N + 1$
CGL nodes reproduces the continuous orthogonality integral exactly.
This discrete orthogonality means that the Chebyshev expansion of a
polynomial of degree $N$ can be computed exactly from its values at
the CGL nodes, using the discrete cosine transform (DCT).

When the function is not a polynomial --- or is a polynomial of degree
higher than $N$ --- the discrete expansion aliases high modes onto low
ones, just as sampling a continuous signal at a finite rate folds
frequencies above the Nyquist limit.  Concretely, the discrete
coefficient $\hat{a}_k$ computed from $N + 1$ samples equals the true
$a_k$ plus contributions from $a_{2N-k}$, $a_{2N+k}$, and so on.
For smooth functions these aliased contributions are negligible
because the high-mode coefficients decay rapidly; for non-smooth
functions, aliasing introduces persistent errors that limit
convergence to algebraic rates.
\xrefnote{The spectral convergence theorem
(\cref{sec:spec-convergence}) quantifies the decay rate of Chebyshev
coefficients as a function of the smoothness of $f$.}

\paragraph{The codebase implementation.}

The \texttt{ChebyshevBasis} struct in \codemodule{spectral}
encapsulates the CGL nodes and the associated differentiation
matrices for $n = N + 1$ collocation points (the struct field
\texttt{n} counts points, not polynomial degree):
%
\begin{lstlisting}[language=Rust]
pub struct ChebyshevBasis {
    pub n: usize,
    pub nodes: Vec<f64>,   // CGL: cos(pi*k/(n-1))
    pub d1: Vec<f64>,      // first derivative matrix
    pub d2: Vec<f64>,      // second derivative matrix
}
\end{lstlisting}
%
The nodes are computed directly from the cosine formula
$x_k = \cos(\pi k / (n-1))$ for $k = 0, \ldots, n-1$, where
$n = N + 1$ is the number of collocation points.  The
differentiation matrices \texttt{d1} and \texttt{d2} --- constructed
using the barycentric formula of Trefethen \cite{Trefethen2000} ---
are the subject of \cref{sec:spec-diffmat}.  The XCTS initial-data
solver uses \texttt{ChebyshevBasis} for elliptic problems on
multi-domain radial grids: each radial subdomain carries its own
set of CGL nodes, and the exponential convergence established above
means that modest polynomial degrees ($N \sim 20$--$30$) suffice to
resolve the metric fields to machine precision.  The Chebyshev
framework thus converts the abstract minimax and orthogonality
results of this section into a concrete computational tool ---
one that underpins every spectral solve in the codebase.
